\subsection{Constrained optimization and reparameterization}
\label{subsec:constrained-optimization}

Many statistical models impose constraints on parameters:
probability vectors live on a simplex, variances must be positive, and covariance matrices must be positive semidefinite.
Abstractly, we consider the constrained optimization problem
\begin{equation}
  \min_{x\in C} f(x),
  \label{eq:constrained-min}
\end{equation}
where $C\subseteq \R^d$ is a constraint set.
Unlike in \eqref{eq:unconstrained-min}, a naive gradient step $x-t\nabla f(x)$ may leave $C$.
Projected methods address this by combining unconstrained updates with a projection back onto the feasible set.
This gives the geometric view of constrained optimization.
We then turn to the algebraic view, where Lagrange multipliers and dual variables encode how the active constraints push back on the objective.

\subsubsection{Projections and projected gradient descent}
\label{subsubsec:projection}
\label{subsubsec:projected-gd}

\begin{definition}[Euclidean projection]
Let $C\subseteq \R^d$ be nonempty, closed, and convex.
The (Euclidean) projection onto $C$ is the map $\Pi_C:\R^d\to C$ defined by
\[
  \Pi_C(z) \coloneqq \argmin_{y\in C}\,\norm{y-z}_2.
\]
\end{definition}

Closedness ensures the minimum is attained; convexity plus the strict convexity of $\norm{\cdot}_2^2$ ensures the minimizer is unique.
Equivalently, $\Pi_C(z)$ is the closest feasible point to $z$ in Euclidean distance.

\begin{algorithm}
  \caption{Projected gradient descent}
  \label{alg:projected-gd}
  \begin{algorithmic}[1]
    \Require Differentiable $f$, closed convex set $C$, initial point $x_0\in C$, step sizes $(t_k)_{k\ge 0}$
    \Ensure Iterate $x_k \in C$
    \For{$k=0,1,2,\dots$ until converged}
      \State $x_{k+1} \gets \Pi_C\bigl(x_k - t_k \nabla f(x_k)\bigr)$
    \EndFor
    \State \Return $x_k$
\end{algorithmic}
\end{algorithm}

Projected gradient descent is conceptually simple: take an unconstrained step, then project back to the feasible set.
In practice, the main difficulty is often computational: computing $\Pi_C$ can be easy for some sets (like boxes) and nontrivial for others.
When projections are expensive, other approaches can be more appropriate.
For example, one can add penalties for constraint violations, use barrier terms (which blow up near the boundary) to keep iterates strictly feasible (see \cref{subsubsec:log-barrier}), or introduce Lagrange multipliers and work with optimality conditions (see \cref{subsubsec:kkt}).
Here we focus on projections because they make the feasibility mechanism explicit.

\begin{figure}[t]
  \centering
  \begin{tikzpicture}[x=1.15cm,y=1.15cm,>=Stealth]
  % Geometry of a projected gradient step on a generic convex set.
  \pgfmathsetmacro{\a}{3.0}
  \pgfmathsetmacro{\b}{2.0}

  % Choose a projection point on a smooth boundary (here: the right-most point of an ellipse),
  % so the projection direction is clearly orthogonal to the tangent.
  \coordinate (xnext) at (\a,0);
  \coordinate (y) at ($ (xnext) + (1.8,0) $);
  \coordinate (xk) at (1.2,-0.85);

  % Tangent segment at x_{k+1} (vertical at the right-most point of an ellipse).
  \coordinate (t1) at ($ (xnext) + (0,1.25) $);
  \coordinate (t2) at ($ (xnext) + (0,-1.25) $);

  % A small right-angle marker between the tangent and the projection direction.
  \pgfmathsetmacro{\s}{0.28}
  \coordinate (ra) at ($ (xnext) + (0,\s) $);
  \coordinate (rb) at ($ (xnext) + (-\s,0) $);
  \coordinate (rc) at ($ (xnext) + (-\s,\s) $);

  % Feasible set C.
  \draw[fill=stat221GrayLight,draw=stat221Gray!65,line width=0.6pt] (0,0) ellipse ({\a} and {\b});
  \node[stat221Gray!70] at (-1.7,1.25) {$C$};

  % Tangent line at the projection point.
  \draw[densely dashed,stat221Gray!70,line width=0.55pt] (t2) -- (t1);

  % Step and projection arrows.
  \draw[stat221Blue,thick,densely dashed,->] (xk) -- (y);
  \draw[stat221Purple,thick,->] (y) -- (xnext);

  % Points.
  \fill[black] (xk) circle (1.45pt);
  \fill[stat221Red] (y) circle (1.45pt);
  \fill[stat221Purple] (xnext) circle (1.55pt);

  % Labels (use a light background to avoid collisions with lines).
  \node[anchor=north east,font=\small,fill=white,fill opacity=0.85,text opacity=1,inner sep=1pt] at (xk) {$x_k$};
  \node[anchor=south west,font=\small,stat221Red,fill=white,fill opacity=0.85,text opacity=1,inner sep=1pt] at (y) {$y$};
  \node[anchor=north east,font=\small,stat221Purple,fill=white,fill opacity=0.85,text opacity=1,inner sep=1pt] at ($ (xnext) + (0.02,-0.02) $) {$x_{k+1}$};

  % Right-angle marker near the projection point.
  \draw[stat221Gray!75,line width=0.55pt] (ra) -- (rc) -- (rb);
\end{tikzpicture}

  \caption{Geometry of a projected gradient step.
  Starting from a feasible iterate $x_k\in C$, we first take an unconstrained step to $y=x_k-t_k\nabla f(x_k)$.
  If $y\notin C$, projected gradient descent maps it back to the feasible set via $x_{k+1}=\Pi_C(y)$, the closest point in $C$ to $y$ in Euclidean distance.
  At the projection point (for smooth boundaries) the correction direction $y-x_{k+1}$ is orthogonal to the boundary.}
  \label{fig:projected-gd-geometry}
\end{figure}

\subsubsection{Affine constraints}
\label{subsubsec:affine-equality-constraints}

A common and especially simple constraint set is an affine set
\[
  C \coloneqq \{x\in\R^d:Ax=b\},
\]
where $A\in\R^{p\times d}$ has full row rank.
In this case the Euclidean projection admits a closed form.

\begin{lemma}[Projection onto an affine set]
\label{lem:proj-affine}
Assume $A$ has full row rank and let $C=\{x\in\R^d:Ax=b\}$.
Then for any $z\in\R^d$,
\[
  \Pi_C(z)=z-A^\top(AA^\top)^{-1}(Az-b).
\]
\end{lemma}
\begin{proof}
The projection solves $\min_x \frac12\norm{x-z}_2^2$ subject to $Ax=b$.
Form the Lagrangian $\mathcal{L}(x,\nu)=\frac12\norm{x-z}_2^2+\nu^\top(Ax-b)$.
Stationarity gives $0=\nabla_x\mathcal{L} = x-z + A^\top \nu$, so $x=z-A^\top\nu$.
Enforcing feasibility yields $Az-AA^\top\nu=b$, hence $\nu=(AA^\top)^{-1}(Az-b)$ since $A$ has full row rank.
\end{proof}

Applying \cref{alg:projected-gd} with this projection keeps all iterates feasible.
Equivalently, the projected gradient step moves in the nullspace $\mathcal{N}(A)$:
\[
  x_{k+1} = x_k - t_k\,\Delta x_k,
  \qquad
  \Delta x_k = \operatorname{Proj}_{\mathcal{N}(A)}\bigl(\nabla f(x_k)\bigr),
\]
so feasibility is preserved for any step size $t_k$.
Here $\mathcal{N}(A)=\{v\in\R^d:Av=0\}$ is the nullspace of $A$, and $\operatorname{Proj}_{\mathcal{N}(A)}$ denotes the Euclidean (orthogonal) projection onto this subspace.
See \citet[\S10.2, Ex.~10.2]{boyd_vandenberghe_2004_convex}.
When $A$ has full row rank, one can write $\operatorname{Proj}_{\mathcal{N}(A)} = I - A^\top(AA^\top)^{-1}A$.

\subsubsection{Simplex and positive-semidefinite projections}
\label{subsubsec:projection-examples}

Each of the projections below is itself a convex optimization problem.
Often, writing a Lagrangian and working out the optimality conditions reveals a simple thresholding structure (an ``active set'' of tight inequalities).
We formalize this viewpoint later via the KKT conditions in \cref{subsubsec:kkt}; here we use it as a computational tool.

A first example is the probability simplex in $\R^d$,
\[
  \Delta \coloneqq \{x\in\R^d:\ x_i\ge 0,\ \textstyle\sum_{i=1}^d x_i = 1\}.
\]
The naive operation ``threshold then renormalize'' is \emph{not} the Euclidean projection in general.
The true projection has a clean structure.

\begin{proposition}[Projection onto the simplex has a threshold form]
\label{prop:simplex-threshold}
Let $z\in\R^d$.
There exists a scalar $\tau\in\R$ such that the Euclidean projection $x^\star=\Pi_\Delta(z)$ satisfies
\begin{equation}
  x^\star_i = \max\{z_i-\tau,0\},
  \qquad i=1,\dots,d,
  \label{eq:simplex-proj-threshold}
\end{equation}
with $\tau$ chosen so that $\sum_i x^\star_i = 1$.
\end{proposition}
\begin{proof}
Consider the convex problem $\min_{x\in\Delta} \frac12\norm{x-z}_2^2$.
Form the Lagrangian with multiplier $\tau$ for the equality constraint and multipliers $\nu_i\ge 0$ for the inequalities:
\[
  \mathcal{L}(x,\tau,\nu)
  = \frac12\norm{x-z}_2^2 + \tau\Bigl(\sum_{i=1}^d x_i - 1\Bigr) - \sum_{i=1}^d \nu_i x_i.
\]
Stationarity gives $0 = \nabla_x \mathcal{L} = x-z + \tau \1 - \nu$, i.e.\ $x_i = z_i-\tau+\nu_i$.
Complementary slackness $\nu_i x_i = 0$ implies:
if $x_i>0$ then $\nu_i=0$ and $x_i=z_i-\tau$; if $x_i=0$ then $\nu_i=\tau-z_i\ge 0$ so $z_i-\tau\le 0$.
Thus $x_i=\max\{z_i-\tau,0\}$.
Finally, feasibility requires $\sum_i x_i=1$, which determines $\tau$.
\end{proof}

\begin{algorithm}
  \caption{Projection onto the simplex $\Delta$ (sorting + threshold)}
  \label{alg:simplex-proj}
  \begin{algorithmic}[1]
    \Require Vector $z\in\R^d$
    \Ensure Projection $x=\Pi_\Delta(z)$
    \State Let $u$ be $z$ sorted in descending order: $u_1\ge \cdots \ge u_d$
    \State $\rho \gets \max\Bigl\{j\in\{1,\dots,d\}:\ u_j - \frac{1}{j}\Bigl(\sum_{i=1}^j u_i - 1\Bigr) > 0\Bigr\}$
    \State $\tau \gets \frac{1}{\rho}\Bigl(\sum_{i=1}^{\rho} u_i - 1\Bigr)$
    \State $x_i \gets \max\{z_i-\tau,0\}$ for $i=1,\dots,d$
    \State \Return $x$
  \end{algorithmic}
\end{algorithm}

\Cref{alg:simplex-proj} implements the threshold form in \cref{prop:simplex-threshold} by finding the active set size $\rho$ and corresponding threshold $\tau$ via sorting.
Complementary slackness identifies an ``active set'' of coordinates that remain positive, and sorting is a convenient way to find its size.
Because $\Delta$ contains interior points (for example $(1/d)\1$), Slater's condition holds for this convex problem, so the KKT conditions in \cref{subsubsec:kkt} indeed characterize the projection.

Another common constraint set is the cone of symmetric positive semidefinite (PSD) matrices $\mathbb{S}^d_+ \subseteq \R^{d\times d}$.
The Euclidean projection (for the Frobenius norm) of a symmetric matrix $Z$ onto $\mathbb{S}^d_+$ is obtained by thresholding eigenvalues.
Here the Frobenius norm is $\norm{M}_F = \left(\sum_{i,j} M_{ij}^2\right)^{1/2}$.
If $Z$ is not symmetric, one typically replaces it by its symmetric part $(Z+Z^\top)/2$ before projecting.

\begin{proposition}[Projection onto the PSD cone]
\label{prop:psd-proj}
Let $Z\in\R^{d\times d}$ be symmetric with eigendecomposition $Z=Q\Lambda Q^\top$, where $\Lambda=\mathrm{diag}(\lambda_1,\dots,\lambda_d)$.
Define $\Lambda_+=\mathrm{diag}(\max\{\lambda_1,0\},\dots,\max\{\lambda_d,0\})$.
Then
\[
  \Pi_{\mathbb{S}^d_+}(Z) = Q\Lambda_+ Q^\top.
\]
\end{proposition}
\begin{proof}
We minimize $\norm{X-Z}_F^2$ over $X\succeq 0$.
By orthogonal invariance of the Frobenius norm,
\[
  \norm{X-Z}_F^2 = \norm{Q^\top X Q - \Lambda}_F^2.
\]
Since $X\succeq 0$ iff $Q^\top X Q \succeq 0$, we can equivalently minimize over $Y\succeq 0$ the quantity $\norm{Y-\Lambda}_F^2$.
For diagonal $\Lambda$, the optimal $Y$ is diagonal as well (off-diagonal entries only increase the Frobenius norm), reducing the problem to projecting $(\lambda_1,\dots,\lambda_d)$ onto $\R_+^d$, i.e.\ replacing negative coordinates by $0$.
\end{proof}

This is the same basic thresholding phenomenon as in \cref{prop:simplex-threshold}, but expressed in an eigenbasis rather than in coordinates.
In the eigen-directions where the PSD constraint is inactive, the projection keeps the eigenvalue unchanged; in the active directions, the eigenvalue is clipped to $0$.

\subsubsection{Reparameterization for strict constraints}
\label{subsubsec:strict-constraints}

The projection $\Pi_C$ is only guaranteed to be well-defined when $C$ is closed.
Strict constraints (e.g.\ $x_i>0$ or $\Sigma\succ 0$) define open sets, which can cause projection-based methods to misbehave near the boundary.
In such cases, a minimizer can fail to exist: the infimum may be approached only at boundary points that are not feasible.
\begin{example}[Open constraints can destroy projections]
If $C=(0,1)\subset\R$ and $z=0$, then $\inf_{x\in C}\norm{x-z}_2=0$ but the infimum is not attained in $C$.
\end{example}
A common alternative is to \emph{reparameterize} the constrained object in terms of unconstrained variables.

A common example is enforcing that a covariance matrix $\Sigma$ is symmetric positive definite.
If $\Sigma\succ 0$ is symmetric, it has a Cholesky factorization $\Sigma=LL^\top$ with $L$ lower triangular and $L_{ii}>0$.
We can enforce strict positivity by writing $L_{ii}=\exp(b_{ii})$ for unconstrained $b_{ii}\in\R$.
Then any choice of parameters produces $\Sigma\succ 0$ automatically.

Similarly, the strict simplex
\[
\Delta^\circ \coloneqq \{x\in\R^d:\ x_i>0,\ \textstyle\sum_{i=1}^d x_i = 1\}
\]
can be reparameterized by the softmax map: for $y\in\R^d$, define
\[
x_i = \frac{e^{y_i}}{\sum_{j=1}^d e^{y_j}},
\qquad i=1,\dots,d.
\]
This enforces $x\in\Delta^\circ$ automatically and turns optimization over $\Delta^\circ$ into unconstrained optimization in $y$ (up to the invariance $y\mapsto y+c\mathbf{1}$).
The next example shows the same idea in a matrix-valued statistical model, where reparameterization removes a fragile positivity constraint and isolates the remaining structural restriction.
\label{subsubsec:gaussian-precision-example}

\begin{example}[Gaussian precision estimation with a sparsity constraint]
Let $x_1,\dots,x_n \in \R^d$ be i.i.d.\ with $x_i\sim \mathcal{N}(\mu,\Sigma)$, where $\Sigma\succ 0$.
Up to an additive constant, the negative log-likelihood is
\begin{equation}
  \mathcal{L}(\mu,\Sigma)
  = \frac{n}{2}\log\det\Sigma + \frac12\sum_{i=1}^n (x_i-\mu)^\top \Sigma^{-1}(x_i-\mu).
  \label{eq:gauss-nll}
\end{equation}
Suppose we impose the (structural) constraint that the $(1,2)$ and $(2,1)$ entries of the precision matrix $\Omega\coloneqq\Sigma^{-1}$ are zero:
\[
  \Omega_{12}=\Omega_{21}=0.
\]
(In Gaussian models, such sparsity constraints correspond to conditional independences.)

Optimizing over $\mu$ first gives $\hat\mu=\bar x$, where $\bar x=\frac1n\sum_{i=1}^n x_i$.
Rewriting the objective in terms of $\Omega=\Sigma^{-1}$ then yields
\begin{equation}
  \mathcal{L}(\bar x,\Omega)
  = \frac{n}{2}\Bigl(-\log\det\Omega + \operatorname{tr}(\Omega S)\Bigr),
  \qquad
  S \coloneqq \frac1n\sum_{i=1}^n (x_i-\bar x)(x_i-\bar x)^\top.
  \label{eq:gauss-nll-omega}
\end{equation}
The matrix $S$ is the sample covariance.
The objective \eqref{eq:gauss-nll-omega} is convex over $\Omega\succ 0$ (since $-\log\det$ is convex and the trace term is linear); see \citet[\S3.1.5]{boyd_vandenberghe_2004_convex}.
Adding the linear constraint $\Omega_{12}=0$ preserves convexity.

To enforce $\Omega\succ 0$, write $\Omega = LL^\top$ with $L$ lower triangular and $L_{ii}=\exp(b_{ii})$.
This is exactly the structure-preserving parameterization from \Cref{subsubsec:cholesky-geometry}, and the corresponding log-determinant and quadratic identities are developed in \Cref{subsubsec:solves-logdets}.
Here the point is that the strict positivity constraint is absorbed into the factor, while the sparsity constraint becomes simpler.
Indeed, because $L$ is lower triangular,
\[
  \Omega_{21} = (LL^\top)_{21} = L_{21}L_{11} = 0 \quad \Rightarrow \quad L_{21}=0,
\]
because $L_{11}=\exp(b_{11})>0$.
So the constrained precision estimation problem reduces to an optimization over the factor entries, with one simple linear constraint in place of the original positivity-plus-sparsity constraints:
\begin{equation}
  \min_{L}\;
  -n\sum_{i=1}^d b_{ii} + \frac{n}{2}\operatorname{tr}(L^\top S L)
  \quad\text{s.t.}\quad
  L \text{ lower triangular},\; L_{ii}=\exp(b_{ii}),\; L_{21}=0.
  \label{eq:gauss-precision-cholesky}
\end{equation}
Concretely, the free variables are the real numbers $b_{11},\dots,b_{dd}$ and the strictly lower-triangular entries of $L$ (with $L_{21}$ fixed to $0$), and every such choice produces a valid precision matrix $\Omega\succ 0$.
This is the main lesson of the example.
Reparameterization removes the fragile positivity constraint, and the remaining structural constraint can be handled either by eliminating the variable or through the KKT framework developed next.
\end{example}

\subsubsection{KKT conditions and equality-constrained Newton}
\label{subsubsec:kkt}
\label{subsubsec:newton-equality-constraints}

The projection examples above already showed the basic multiplier idea.
If we want to minimize \(f(x)\) subject to one equality constraint \(h(x)=0\), we form the Lagrangian
\[
  \mathcal{L}(x,\nu)=f(x)+\nu h(x).
\]
At a constrained optimum, the gradient of \(f\) need not vanish.
Instead it is balanced by the gradient of the constraint:
\[
  \nabla f(x^\star)+\nu^\star \nabla h(x^\star)=0.
\]
The multiplier records how strongly the constraint pushes back against further decrease of the objective.
The KKT conditions are the general version of this first-order balance.

Consider a problem with inequality and equality constraints,
\begin{equation}
  \min_{x\in\R^d} f(x)
  \quad\text{s.t.}\quad
  g_i(x)\le 0,\ i=1,\dots,m,
  \qquad
  h_j(x)=0,\ j=1,\dots,p,
  \label{eq:ineq-eq-constrained}
\end{equation}
where $f,g_i,h_j$ are differentiable.
The constrained problem \eqref{eq:ineq-eq-constrained} is called the \emph{primal problem}, and \(x\) is the \emph{primal variable}.
For multipliers $\lambda\in\R_+^m$ (i.e.\ $\lambda_i\ge 0$) and $\nu\in\R^p$, define the Lagrangian
\[
  \mathcal{L}(x,\lambda,\nu)
  \coloneqq f(x) + \sum_{i=1}^m \lambda_i g_i(x) + \sum_{j=1}^p \nu_j h_j(x).
\]
For fixed multipliers, \(\mathcal{L}(x,\lambda,\nu)\) is just an unconstrained objective in \(x\): the extra terms encode how much we penalize or reward constraint violation.
The same construction also defines the \emph{dual function}
\[
  q(\lambda,\nu)\coloneqq \inf_{x\in\R^d}\,\mathcal{L}(x,\lambda,\nu),
\]
and the associated \emph{dual problem}
\[
  \max_{\lambda\in\R_+^m,\ \nu\in\R^p}\; q(\lambda,\nu).
\]
The dual function is the best lower bound on the primal objective that the multipliers can certify, and the dual problem searches for the strongest such bound.
For each fixed $(\lambda,\nu)$, the dual function is obtained by minimizing the Lagrangian over the primal variable \(x\).
This automatically makes $q$ a concave function of $(\lambda,\nu)$, even if the primal problem itself is not convex.
Let $p^\star$ denote the optimal value of \eqref{eq:ineq-eq-constrained}, and let $d^\star$ denote the optimal value of the dual.
Weak duality gives $d^\star\le p^\star$.
Indeed, if $x$ is primal feasible and $\lambda\ge 0$, then
\[
  \mathcal{L}(x,\lambda,\nu)
  = f(x) + \sum_{i=1}^m \lambda_i g_i(x) + \sum_{j=1}^p \nu_j h_j(x)
  \le f(x),
\]
because the inequality-constraint terms are nonpositive and the equality-constraint terms vanish.
Taking the infimum over $x$ gives $q(\lambda,\nu)\le f(x)$.
More generally, for any primal feasible $x$ and dual feasible $(\lambda,\nu)$ (meaning $\lambda\ge 0$ and $q(\lambda,\nu)>-\infty$) we have $q(\lambda,\nu)\le f(x)$, so the nonnegative quantity $f(x)-q(\lambda,\nu)$ is a computable \emph{duality gap} certificate; see \citet[\S5.1.2--\S5.1.3]{boyd_vandenberghe_2004_convex}.

\begin{example}[Dual problem for affine projection]
\label{ex:affine-projection-dual}
Consider the affine projection problem from \cref{lem:proj-affine},
\[
  \min_{x\in\R^d} \frac12\norm{x-z}_2^2
  \qquad\text{s.t.}\qquad
  Ax=b,
\]
where $A$ has full row rank.
Its Lagrangian is $\mathcal{L}(x,\nu)=\frac12\norm{x-z}_2^2+\nu^\top(Ax-b)$.
Minimizing over $x$ gives $x=z-A^\top\nu$, and substituting this back yields the dual function
\[
  q(\nu)=\nu^\top(Az-b)-\frac12\nu^\top AA^\top \nu.
\]
The dual problem is therefore the unconstrained concave quadratic
\[
  \max_{\nu\in\R^p}\;
  \nu^\top(Az-b)-\frac12\nu^\top AA^\top \nu.
\]
Its maximizer is $\nu^\star=(AA^\top)^{-1}(Az-b)$, and substituting back into $x=z-A^\top\nu$ recovers the projection formula from \cref{lem:proj-affine}.
This example makes the roles clear: the primal variable is the candidate point $x$, while the dual variable $\nu$ enforces the feasibility condition $Ax=b$.
\end{example}

The Karush--Kuhn--Tucker (KKT) conditions are the standard first-order optimality conditions associated with \eqref{eq:ineq-eq-constrained}.
For convex problems they can give an exact, checkable characterization of optimality; for nonconvex problems they should be viewed as necessary conditions (under mild regularity of the constraints), not as a guarantee of global optimality.
We say that $(x^\star,\lambda^\star,\nu^\star)$ satisfies the KKT conditions if
\begin{equation}
  \begin{aligned}
    g_i(x^\star) &\le 0, && i=1,\dots,m, \\
    h_j(x^\star) &= 0, && j=1,\dots,p, \\
    \lambda_i^\star &\ge 0, && i=1,\dots,m, \\
    \lambda_i^\star g_i(x^\star) &= 0, && i=1,\dots,m, \\
    \nabla f(x^\star) + \sum_{i=1}^m \lambda_i^\star \nabla g_i(x^\star) + \sum_{j=1}^p \nu_j^\star \nabla h_j(x^\star) &= 0.
  \end{aligned}
  \label{eq:kkt-conditions}
\end{equation}

The last line in \eqref{eq:kkt-conditions} is called \emph{stationarity}: at optimality, the objective gradient is balanced by a linear combination of constraint gradients.
Complementary slackness, $\lambda_i^\star g_i(x^\star)=0$, captures the ``active set'' idea:
if constraint $g_i(x)\le 0$ is \emph{inactive} at $x^\star$ (meaning $g_i(x^\star)<0$), then necessarily $\lambda_i^\star=0$, so it plays no role in stationarity.
Only constraints that are tight ($g_i(x^\star)=0$) can have nonzero multipliers; we call such constraints \emph{active}.
These multipliers appear in the stationarity equation as a linear combination of the gradients of the active constraints.
For a differentiable constraint $g_i$, the vector $\nabla g_i(x^\star)$ is orthogonal to the boundary $\{x:\ g_i(x)=0\}$ at $x^\star$, so it acts as a (local) normal vector to the constraint surface.
For an inequality constraint $g_i(x)\le 0$, $\nabla g_i(x^\star)$ points locally in the direction where $g_i$ increases, i.e.\ toward constraint violation.

This matches the basic geometry of constrained minima.
If $x^\star$ lies in the interior of the feasible region, then all inequalities are inactive, complementary slackness forces $\lambda^\star=0$, and stationarity reduces to $\nabla f(x^\star)=0$.
If $x^\star$ lies on the boundary, then stationarity says that $-\nabla f(x^\star)$ can be written as a linear combination of constraint gradients, with nonnegative coefficients on the active inequalities.
In other words, any direction that decreases $f$ must point (at least a little) out of the feasible set, and the multipliers encode how the active constraints ``push back'' against the objective gradient.
In the unconstrained case ($m=p=0$), \eqref{eq:kkt-conditions} reduces to $\nabla f(x^\star)=0$.
See \cref{fig:kkt-geometry} for a two-dimensional picture of this balance.

\begin{figure}[H]
  \centering
  \begin{tikzpicture}[x=1.1cm,y=1.1cm,>=Stealth]
  % Two-panel sketch of the geometric content of KKT stationarity.

	  % --- (a) Smooth boundary, single active inequality ---
	  \begin{scope}[shift={(0,0)}]
	    \pgfmathsetmacro{\a}{2.35}
	    \pgfmathsetmacro{\b}{1.55}

    \coordinate (xstar) at (\a,0);
    \coordinate (t1) at ($ (xstar) + (0,1.2) $);
    \coordinate (t2) at ($ (xstar) + (0,-1.2) $);

    % Feasible set.
    \draw[fill=stat221GrayLight,draw=stat221Gray!65,line width=0.6pt] (0,0) ellipse ({\a} and {\b});
    \node[stat221Gray!70] at (-1.55,1.1) {$C$};

    % Tangent line at x^* (smooth boundary).
    \draw[densely dashed,stat221Gray!70,line width=0.55pt] (t2) -- (t1);

	    % Point x^*.
	    \fill[black] (xstar) circle (1.45pt);
	    \node[anchor=south east,font=\small,fill=white,fill opacity=0.9,text opacity=1,inner sep=1pt] at (xstar) {$x^\star$};

	    % Single visual statement of stationarity: -∇f aligns with the outward normal ∇g.
	    \draw[stat221Blue,thick,->] (xstar) -- ++(1.25,0)
	      node[pos=0.62,anchor=south,font=\small,stat221Blue,fill=white,fill opacity=0.9,text opacity=1,inner sep=1pt]
	      {$-\nabla f(x^\star)$};
	    \draw[stat221Purple,thick,->] (xstar) -- ++(0.9,0)
	      node[pos=0.62,anchor=north,font=\small,stat221Purple,fill=white,fill opacity=0.9,text opacity=1,inner sep=1pt]
	      {$\lambda^\star\nabla g(x^\star)$};

	    % Right-angle marker between the tangent and the normal direction.
	    \pgfmathsetmacro{\s}{0.28}
	    \coordinate (ra) at ($ (xstar) + (0,\s) $);
	    \coordinate (rb) at ($ (xstar) + (-\s,0) $);
	    \coordinate (rc) at ($ (xstar) + (-\s,\s) $);
	    \draw[stat221Gray!75,line width=0.55pt] (ra) -- (rc) -- (rb);

	    \node[anchor=north west,font=\small] at (-2.45,1.9) {(a)};
	  \end{scope}

  % --- (b) Corner, two active inequalities ---
	  \begin{scope}[shift={(6.6,0)}]
	    \pgfmathsetmacro{\R}{2.25}
	    \pgfmathsetmacro{\Rc}{1.65}
	    \coordinate (xstar) at (0,0);

	    % Feasible set with a corner at x^* (intersection of two constraints): truncated by an arc.
	    \path[fill=stat221GrayLight] (0,0) -- (\R,0) arc (0:90:\R) -- cycle;
	    \draw[stat221Gray!65,line width=0.6pt] (0,0) -- (\R,0);
	    \draw[stat221Gray!65,line width=0.6pt] (0,0) -- (0,\R);
	    \draw[stat221Gray!35,line width=0.55pt] (\R,0) arc (0:90:\R);
	    \node[stat221Gray!70] at (1.05,1.15) {$C$};

    % Point x^*.
    \fill[black] (xstar) circle (1.45pt);
    \node[anchor=south east,font=\small,fill=white,fill opacity=0.9,text opacity=1,inner sep=1pt] at (xstar) {$x^\star$};

		    % Shading for the set of nonnegative combinations of the active normals.
		    \path[fill=stat221Purple!15] (0,0) -- (-\Rc,0) -- (0,-\Rc) -- cycle;
		    \draw[stat221Purple!35,line width=0.45pt] (0,0) -- (-\Rc,0);
		    \draw[stat221Purple!35,line width=0.45pt] (0,0) -- (0,-\Rc);

	    % Active constraint gradients (outward normals) at the corner.
	    \draw[stat221Purple,thick,->] (xstar) -- ++(-1.45,0)
	      node[anchor=east,font=\small,stat221Purple,fill=white,fill opacity=0.9,text opacity=1,inner sep=1pt] {$\nabla g_1(x^\star)$};
    \draw[stat221Purple,thick,->] (xstar) -- ++(0,-1.45)
      node[anchor=north,font=\small,stat221Purple,fill=white,fill opacity=0.9,text opacity=1,inner sep=1pt] {$\nabla g_2(x^\star)$};

	    % Stationarity: -∇f lies in the normal cone spanned by the active normals.
	    \draw[stat221Blue,thick,->] (xstar) -- ++(-1.05,-0.85)
	      node[anchor=north east,font=\small,stat221Blue,fill=white,fill opacity=0.9,text opacity=1,inner sep=1pt] {$-\nabla f(x^\star)$};

    \node[anchor=north west,font=\small] at (-2.05,1.9) {(b)};
  \end{scope}
\end{tikzpicture}

  \caption{Geometry of KKT stationarity and complementary slackness.
  Left: with a single active inequality at a smooth boundary point $x^\star$, stationarity forces $-\nabla f(x^\star)$ to align with the constraint normal $\nabla g(x^\star)$, with some scale $\lambda^\star\ge 0$.
  Right: at a corner with two active inequalities, stationarity forces $-\nabla f(x^\star)$ to lie in the set of nonnegative combinations of the active normals $\nabla g_1(x^\star)$ and $\nabla g_2(x^\star)$.
  This same ``normal direction'' geometry underlies Euclidean projection steps, such as the projected update in \cref{fig:projected-gd-geometry}.}
  \label{fig:kkt-geometry}
\end{figure}

This same KKT geometry also explains the geometry of Euclidean projection.
The projection $\Pi_C(z)$ is defined as the minimizer of $\min_{x\in C} \frac12\norm{x-z}_2^2$, which is a convex problem when $C$ is convex.
When $C$ is described by differentiable constraints, applying the KKT stationarity equation to this projection problem shows that $z-\Pi_C(z)$ must be a linear combination of gradients of the active constraints at the projection point.
For example, if $C=\{x:\ g(x)\le 0\}$ with differentiable $g$ and $x^\star=\Pi_C(z)$ lies on the boundary $g(x^\star)=0$, then there exists $\lambda^\star\ge 0$ such that
\[
  x^\star - z + \lambda^\star \nabla g(x^\star)=0,
  \qquad\text{so}\qquad
  z-x^\star = \lambda^\star \nabla g(x^\star).
\]
This is exactly the normal-direction geometry illustrated in \cref{fig:projected-gd-geometry}: for a smooth boundary, the projection correction is orthogonal to the tangent.

\begin{theorem}[KKT optimality conditions for convex problems]
\label{thm:kkt-convex}
Consider the constrained problem \eqref{eq:ineq-eq-constrained}.
Assume that $f$ and $g_i$ are convex, that each equality constraint $h_j$ is affine (linear plus a constant), and that there exists a \emph{strictly feasible} point $\tilde x$ with $g_i(\tilde x)<0$ for all $i$ and $h_j(\tilde x)=0$ for all $j$ (Slater's condition).
Then $x^\star$ solves \eqref{eq:ineq-eq-constrained} if and only if there exist multipliers $\lambda^\star\in\R_+^m$ and $\nu^\star\in\R^p$ such that $(x^\star,\lambda^\star,\nu^\star)$ satisfies the KKT conditions \eqref{eq:kkt-conditions}.
See \citet[\S5.5.3]{boyd_vandenberghe_2004_convex}.
\end{theorem}

Under Slater's condition, strong duality holds ($p^\star=d^\star$) and the KKT multipliers can be viewed as a dual certificate of optimality (zero duality gap); see \citet[\S5.5.3]{boyd_vandenberghe_2004_convex}.
This is the basic \emph{primal--dual} viewpoint:
one seeks a primal variable $x$ and dual multipliers $(\lambda,\nu)$ that simultaneously satisfy primal feasibility, dual feasibility, complementary slackness, and stationarity.
The KKT system is precisely the set of equations and inequalities that couples these two sides of the problem.

The KKT conditions are not only a certificate of optimality; they can also be viewed as a system of equations to solve.
A particularly transparent case is the equality constrained problem
\begin{equation}
  \min_{x\in\R^d} f(x)
  \quad\text{s.t.}\quad
  Ax=b,
  \label{eq:equality-constrained}
\end{equation}
where $A\in\R^{p\times d}$ has full row rank $p<d$ and $f$ is twice differentiable.
Projected gradient descent can enforce feasibility via the closed-form projection in \cref{lem:proj-affine}.

Newton's method applies the same feasibility idea to the local quadratic model of $f$.
Assume $x$ is feasible ($Ax=b$) and let $H=\Hess f(x)$.
The equality constrained Newton step $\Delta x_{\mathrm{nt}}$ is the minimizer of the quadratic model subject to preserving feasibility to first order, i.e.\ staying in the tangent space:
\[
  \Delta x_{\mathrm{nt}}
  \in \argmin_{\Delta x}\; \nabla f(x)^\top \Delta x + \frac12 \Delta x^\top H \Delta x
  \quad\text{s.t.}\quad
  A\Delta x = 0.
\]
Introducing a Lagrange multiplier $\nu\in\R^p$ for the constraint $A\Delta x=0$, the KKT conditions (see \cref{subsubsec:kkt}) yield the linear system
\begin{equation}
  \begin{bmatrix}
    H & A^\top \\
    A & 0
  \end{bmatrix}
  \begin{bmatrix}
    \Delta x_{\mathrm{nt}} \\
    \nu
  \end{bmatrix}
  =
  \begin{bmatrix}
    -\nabla f(x) \\
    0
  \end{bmatrix}.
  \label{eq:equality-newton-kkt}
\end{equation}
The top block row is the stationarity condition for the quadratic model (with the equality constraint incorporated via $A^\top\nu$), while the bottom block row enforces feasibility $A\Delta x_{\mathrm{nt}}=0$.
The block matrix in \eqref{eq:equality-newton-kkt} is called the (equality-constrained) KKT matrix.
If $A$ has full row rank and $H$ is positive definite on $\mathcal{N}(A)$ (i.e.\ $\Delta x\ne 0$ and $A\Delta x=0$ implies $\Delta x^\top H \Delta x>0$), then this matrix is nonsingular and \eqref{eq:equality-newton-kkt} has a unique solution; see \citet[\S10.2]{boyd_vandenberghe_2004_convex}.
One typically combines $\Delta x_{\mathrm{nt}}$ with a backtracking line search on $f$ (feasibility is preserved since $A\Delta x_{\mathrm{nt}}=0$).

\subsubsection{Barrier methods}
\label{subsubsec:log-barrier}

For inequality constraints, another common approach is to keep the iterates strictly inside the feasible region by adding a term that blows up at the boundary.
One exact but nonsmooth way to express this is to rewrite each inequality $g_i(x)\le 0$ with the indicator of the negative half-line,
\[
  I_{(-\infty,0]}(u)
  \coloneqq
  \begin{cases}
    0, & u\le 0,\\
    +\infty, & u>0.
  \end{cases}
\]
Then the constrained problem can be written equivalently as minimizing
\[
  f(x) + \sum_{i=1}^m I_{(-\infty,0]}\bigl(g_i(x)\bigr)
  \qquad \text{subject to} \qquad Ax=b.
\]
This makes the role of the inequalities completely explicit, but it is not suitable for Newton-like updates because the indicator is flat on the interior and singular at the boundary.
The log barrier replaces this hard wall by the smooth surrogate $-(1/t)\log(-u)$, which is finite on the interior and diverges as $u\uparrow 0$.
Multiplying the resulting surrogate objective by $t$ leads to the standard barrier form $t f(x)+\phi(x)$.
This is the standard starting point for barrier methods; see \citet[\S11.3]{boyd_vandenberghe_2004_convex}.
The model example is the positive half-line \(x>0\), where the function \(-\log x\) is mild in the interior and diverges as \(x\downarrow 0\).
Barrier methods use the same idea in higher-dimensional constrained problems.
To make the discussion concrete, consider the convex problem
\begin{equation}
  \min_{x\in\R^d} f(x)
  \quad\text{s.t.}\quad
  g_i(x)\le 0,\ i=1,\dots,m,
  \qquad
  Ax=b,
  \label{eq:barrier-primal}
\end{equation}
where $f$ and the $g_i$ are convex and twice differentiable, and $A$ has full row rank.
We assume there exists a strictly feasible point $\tilde x$ with $g_i(\tilde x)<0$ for all $i$ and $A\tilde x=b$.

For the inequalities $g_i(x)\le 0$, the logarithmic barrier is
\[
  \phi(x) \coloneqq -\sum_{i=1}^m \log\bigl(-g_i(x)\bigr),
  \qquad
  \operatorname{dom}\phi = \{x:\ g_i(x)<0\ \text{for all}\ i\}.
\]
For a parameter $t>0$, define the barrier objective
\[
  F_t(x)\coloneqq t f(x)+\phi(x).
\]
The associated centering problem is
\begin{equation}
  x^\star(t)\in \argmin_{x}\; F_t(x)
  \quad\text{s.t.}\quad
  Ax=b.
  \label{eq:centering-problem}
\end{equation}
The barrier term $\phi$ keeps iterates strictly feasible, while increasing $t$ pushes the minimizer of $F_t$ toward a solution of the original constrained problem.

For later use, it is helpful to record the derivatives of $\phi$:
\begin{align*}
  \nabla \phi(x)
  &= -\sum_{i=1}^m \frac{1}{g_i(x)}\,\nabla g_i(x), \\
  \nabla^2 \phi(x)
  &= \sum_{i=1}^m\left(\frac{1}{g_i(x)^2}\,\nabla g_i(x)\nabla g_i(x)^\top - \frac{1}{g_i(x)}\,\nabla^2 g_i(x)\right).
\end{align*}
Because $g_i(x)<0$, the coefficients $-1/g_i(x)$ are positive, so (under convexity of $g_i$) the barrier Hessian is positive semidefinite.

The minimizer $x^\star(t)$ of the centering problem traces out the \emph{central path} as $t$ varies.
It also yields a dual feasible point and a simple duality-gap certificate.
This central-path viewpoint and the resulting gap bound are the core conclusions of \citet[\S11.3]{boyd_vandenberghe_2004_convex}.
Let $p^\star$ denote the optimal value of \eqref{eq:barrier-primal}.
The KKT conditions for \eqref{eq:centering-problem} state that there exists a multiplier $\hat\nu(t)\in\R^p$ such that
\begin{equation}
  t\nabla f\bigl(x^\star(t)\bigr) + \nabla \phi\bigl(x^\star(t)\bigr) + A^\top \hat\nu(t)=0,
  \qquad
  Ax^\star(t)=b.
  \label{eq:centering-kkt}
\end{equation}
Using the explicit form of $\nabla \phi$, we can rewrite stationarity as
\[
  \nabla f\bigl(x^\star(t)\bigr) + \sum_{i=1}^m \lambda_i(t)\,\nabla g_i\bigl(x^\star(t)\bigr) + A^\top \nu(t)=0,
\]
where we define
\[
  \lambda_i(t)\coloneqq -\frac{1}{t\,g_i(x^\star(t))}>0,
  \qquad
  \nu(t)\coloneqq \frac{1}{t}\hat\nu(t).
\]
Consider the Lagrangian for \eqref{eq:barrier-primal},
\[
  \mathcal{L}(x,\lambda,\nu)\coloneqq f(x) + \sum_{i=1}^m \lambda_i g_i(x) + \nu^\top(Ax-b).
\]
In the convex setting, $\lambda(t)\ge 0$ implies that $\mathcal{L}(x,\lambda(t),\nu(t))$ is convex in $x$.
The stationarity equation above is exactly $\nabla_x \mathcal{L}(x^\star(t),\lambda(t),\nu(t))=0$.
Since $\mathcal{L}(\cdot,\lambda(t),\nu(t))$ is convex and differentiable, this implies that $x^\star(t)$ minimizes $\mathcal{L}(\cdot,\lambda(t),\nu(t))$, so the dual function value $q(\lambda(t),\nu(t))$ from \cref{subsubsec:kkt} is attained at $x^\star(t)$.
Since $Ax^\star(t)=b$, this gives
\[
  q\bigl(\lambda(t),\nu(t)\bigr)=f\bigl(x^\star(t)\bigr)+\sum_{i=1}^m \lambda_i(t)\,g_i\bigl(x^\star(t)\bigr).
\]
Moreover,
\[
  \sum_{i=1}^m \lambda_i(t)\,g_i\bigl(x^\star(t)\bigr)
  = \sum_{i=1}^m \left(-\frac{1}{t\,g_i(x^\star(t))}\right) g_i\bigl(x^\star(t)\bigr)
  = -\frac{m}{t},
\]
so the duality gap for the primal--dual pair $\bigl(x^\star(t),\lambda(t),\nu(t)\bigr)$ equals $m/t$.
In particular, $f(x^\star(t)) - p^\star \le m/t$.
Thus, in principle, computing a central point with $m/t\le \varepsilon$ yields an $\varepsilon$-suboptimal solution of \eqref{eq:barrier-primal}.
See \citet[\S11.2]{boyd_vandenberghe_2004_convex}.

The barrier method computes approximate central points for an increasing sequence of $t$ values, terminating once the duality-gap bound $m/t$ is below a tolerance.
\begin{algorithm}
  \caption{Barrier method (logarithmic barrier)}
  \label{alg:barrier-method}
  \begin{algorithmic}[1]
    \Require Convex $f$ and $g_i$, affine constraint $Ax=b$, strictly feasible $x_0$ with $g_i(x_0)<0$
    \Require Parameters $t_0>0$, $\mu>1$, tolerance $\varepsilon>0$
    \State $x\gets x_0$, $t\gets t_0$
    \While{$m/t > \varepsilon$}
      \State Approximately solve \eqref{eq:centering-problem} starting at $x$ (centering step), and update $x$ to the computed solution
      \State $t \gets \mu t$
    \EndWhile
    \State \Return $x$
  \end{algorithmic}
\end{algorithm}

Within each centering step, one can apply Newton's method to minimize $F_t$ subject to $Ax=b$.
At a current strictly feasible point $x$ with $Ax=b$, let $r\coloneqq \nabla F_t(x)$ and $H\coloneqq \nabla^2 F_t(x)$.
The equality-constrained Newton step solves the KKT linear system
\begin{equation}
  \begin{bmatrix}
    H & A^\top \\
    A & 0
  \end{bmatrix}
  \begin{bmatrix}
    \Delta x \\
    w
  \end{bmatrix}
  =
  \begin{bmatrix}
    -r \\
    0
  \end{bmatrix}.
  \label{eq:barrier-centering-newton}
\end{equation}
One then updates $x\gets x+\alpha \Delta x$ with a backtracking line search that (i) preserves strict feasibility $g_i(x+\alpha\Delta x)<0$ and (ii) decreases $F_t$ (e.g.\ using an Armijo condition).
See \citet[\S11.3]{boyd_vandenberghe_2004_convex}.
